\subsection{Definition of the Fourier transform}

DEFINITION 3. --- Let \(\mu \in \mathcal{M}^{1}(\mathrm{G})\) be a bounded complex measure on G. The Fourier transform of \(\mu\) is the function \(\mathcal{F}_{\mathrm{G}}(\mu)\) on \(\widehat{\mathrm{G}}\) defined by

\[
\mathcal {F} _ {\mathrm{G}} (\mu) (\widehat {x}) = \int_ {\mathrm{G}} \overline {{\langle \widehat {x} , x \rangle}} d \mu (x).\tag{5}
\]

We call the Fourier cotransform of \(\mu\) the function \(\overline{\mathcal{F}}_{\mathrm{G}}(\mu)\) on \(\widehat{\mathrm{G}}\) defined by

\[
\overline {{\mathcal {F}}} _ {\mathrm{G}} (\mu) (\widehat {x}) = \int_ {\mathrm{G}} \langle \widehat {x}, x \rangle d \mu (x).\tag{6}
\]

When there is no ambiguity regarding the group G under consideration, we will also write \(\mathcal{F}(\mu)\) and \(\overline{\mathcal{F}}(\mu)\). We also sometimes denote by \(\widehat{\mu} = \mathcal{F}_{\mathrm{G}}(\mu)\).

PROPOSITION 3. --- For every measure \(\mu \in \mathcal{M}^{1}(\mathrm{G})\), the functions \(\mathcal{F}_{\mathrm{G}}(\mu)\) and \(\overline{\mathcal{F}}_{\mathrm{G}}(\mu)\) are continuous and bounded. The maps \(\mathcal{F}_{\mathrm{G}}: \mu \mapsto \mathcal{F}_{\mathrm{G}}(\mu)\) and \(\overline{\mathcal{F}}_{\mathrm{G}}: \mu \mapsto \overline{\mathcal{F}}_{\mathrm{G}}(\mu)\) are continuous morphisms of the involutive algebra \(\mathcal{M}^{1}(\mathrm{G})\) into the involutive algebra of bounded continuous functions on \(\widehat{\mathrm{G}}\) (example 1 of I, p.~99).

Let \(\mu \in \mathcal{M}^1 (\mathrm{G})\). For every \(\chi \in \widehat{\mathrm{G}}\), one has

\[
| \mathcal {F} (\mu) (\chi) | = \left| \int_ {\mathrm{G}} \overline {{\langle \chi , x \rangle}} d \mu (x) \right| \leqslant \| \mu \| _ {1},\tag{7}
\]

so the Fourier transform of \(\mu\) is bounded. Similarly, one verifies that \(\overline{\mathcal{F}}(\mu)\) is bounded.

If \(\chi\) tends to \(\chi_0\) in \(\widehat{\mathbf{G}}\), the function \(\chi\) on \(\mathbf{G}\) tends to \(\chi_0\) uniformly on every compact set while remaining bounded by the constant function 1, which belongs to \(\mathrm{L}^1 (\mathrm{G},\mu)\). By Lebesgue's theorem (INT, IV, §3, n°7, th. 6), it follows that \(\mathcal{F}(\mu)(\chi)\) tends to \(\mathcal{F}(\mu)(\chi_0)\). Therefore \(\mathcal{F}(\mu)\) is continuous. The same is true of \(\overline{\mathcal{F}} (\mu)\).

For every \(\chi \in \widehat{\mathbf{G}}\), the map \(\mu \mapsto \chi(\mu) = \int \langle \chi, x \rangle d\mu(x)\) is a Hermitian character of \(\mathcal{M}^1(\mathbf{G})\) (lemma 1 of II, p.~202). This implies that \(\mathcal{F}\) and \(\overline{\mathcal{F}}\) are involutive algebra morphisms of \(\mathcal{M}^1(\mathbf{G})\) into the involutive algebra of bounded continuous functions on \(\widehat{\mathbf{G}}\). The inequality (7) shows that these morphisms are continuous.

The Fourier transform of G (resp. the Fourier cotransform of G) is the mapping \(\mu \mapsto \mathcal{F}_{\mathrm{G}}(\mu)\) (resp. the map \(\mu \mapsto \overline{\mathcal{F}}_{\mathrm{G}}(\mu)\)) from \(\mathcal{M}^1 (\mathrm{G})\) into \(\mathcal{C}_b(\widehat{\mathrm{G}})\).

Let us note some useful formulas for \(\mu\in\mathcal{M}^{1}(G)\), \(x\in G\), and \(\chi\in\widehat{G}\):

\begin{enumerate}
\def\labelenumi{(\arabic{enumi})}
\setcounter{enumi}{7}
\tightlist
\item
\end{enumerate}

\[
\overline {{\mathcal {F}}} (\mu) (\chi) = \mathcal {F} (\mu) (\chi^ {- 1}) = \overline {{\mathcal {F} (\overline {{\mu}}) (\chi)}},\tag{9}
\]

\[
\| \mathcal {F} (\mu) \| _ {\infty} = \| \overline {{\mathcal {F}}} (\mu) \| _ {\infty} \leqslant \| \mu \| _ {1},\tag{10}
\]

\[
\left\{ \begin{array}{l} \mathcal {F} (\varepsilon_ {x}) (\chi) = \overline {{\langle \chi , x \rangle}}, \\ \overline {{\mathcal {F}}} (\varepsilon_ {x}) (\chi) = \langle \chi , x \rangle , \end{array} \right.
\]

(in particular \(\mathcal{F}(\varepsilon_e) = \overline{\mathcal{F}} (\varepsilon_e) = 1\) ),

\begin{enumerate}
\def\labelenumi{(\arabic{enumi})}
\setcounter{enumi}{10}
\tightlist
\item
\end{enumerate}

\[
\left\{ \begin{array}{l} \mathcal {F} (\varepsilon_ {x} * \mu) (\chi) = \overline {{\langle \chi , x \rangle}} \mathcal {F} (\mu) (\chi), \\ \overline {{\mathcal {F}}} (\varepsilon_ {x} * \mu) (\chi) = \langle \chi , x \rangle \overline {{\mathcal {F}}} (\mu) (\chi), \end{array} \right.\tag{12}
\]

\[
\left\{ \begin{array}{l} \mathcal {F} (\chi \cdot \mu) = \varepsilon_ {\chi} * \mathcal {F} (\mu), \\ \overline {{\mathcal {F}}} (\chi \cdot \mu) = \varepsilon_ {\chi^ {- 1}} * \overline {{\mathcal {F}}} (\mu). \end{array} \right.
\]

Formulas (8), (9), (10), and (11) follow from the definitions. Let us prove the first of formulas (12), the second being analogous. For all \(\xi\) in \(\widehat{G}\), the equalities

\[
\begin{array}{c} \mathcal {F} (\chi \cdot \mu) (\xi) = \int_ {\mathrm{G}} \overline {{\langle \xi , x \rangle}} \langle \chi , x \rangle d \mu (x) = \int_ {\mathrm{G}} \overline {{\langle \xi \chi^ {- 1} , x \rangle}} d \mu (x) \\ = \mathcal {F} (\mu) (\xi \chi^ {- 1}) = (\varepsilon_ {\chi} * \mathcal {F} (\mu)) (\xi). \end{array}
\]

Note moreover that for all \(\chi\in\widehat{\mathrm{G}}\) and all measures \(\mu\) and \(\nu\) in \(\mathcal{M}^{1}(\mathrm{G})\), one has

\[
(\varepsilon_ {\chi} * \mathcal {F} (\mu)) (\varepsilon_ {\chi} * \mathcal {F} (\nu)) = \varepsilon_ {\chi} * (\mathcal {F} (\mu) \mathcal {F} (\nu)),\tag{13}
\]

since both sides of this equality are functions on \(\widehat{G}\) whose value at \(\xi\in\widehat{G}\) is

\[
\mathcal {F} (\mu) (\xi \chi^ {- 1}) \mathcal {F} (\nu) (\xi \chi^ {- 1}).
\]

Let H be a locally compact abelian group, and let \(\varphi\colon\mathrm{G}\to\mathrm{H}\) be a continuous morphism. Let \(\mu\in\mathcal{M}^{1}(\mathrm{G})\). The image measure \(\varphi(\mu)\) is defined (INT, V, §6, no. 1, remark 1), and it follows that \(\mathcal{F}_{\mathrm{H}}(\varphi(\mu))=\mathcal{F}_{\mathrm{G}}(\mu)\circ\widehat{\varphi}\) (cf. INT, V, §6, no. 4, prop. 7).

By restriction to the subalgebra \(L^{1}(G)\) of \(\mathcal{M}^{1}(G)\), we obtain the definitions of the Fourier transform and the Fourier cotransform on \(L^{1}(G)\). We therefore have, for \(f \in L^{1}(G)\) and \(\chi \in \widehat{G}\):

\[
\mathcal {F} _ {\mathrm{G}} (f) (\chi) = \int_ {\mathrm{G}} \overline {{\langle \chi , x \rangle}} f (x) d x, \quad \overline {{\mathcal {F}}} _ {\mathrm{G}} (f) (\chi) = \int_ {\mathrm{G}} \langle \chi , x \rangle f (x) d x.\tag{14}
\]

In particular, \(\mathcal{F}_{\mathrm{G}}(f) = \overline{\overline{\mathcal{F}}_{\mathrm{G}}(\overline{f})}\). We also have

\[
\overline {{\mathcal {F}}} (f) (\chi) = \chi (f)\tag{15}
\]

for all \(f \in \mathrm{L}^{1}(\mathrm{G})\) and every \(\chi \in \widehat{G}\).

Let \(\sigma\) be an automorphism of G and \(\Delta\) the modulus of \(\sigma\) (INT, VII, §1, no. 4, def. 4). For \(f \in \mathrm{L}^1(\mathrm{G})\), one has

\[
\mathcal {F} (f \circ \sigma) = \Delta^ {- 1} \mathcal {F} (f) \circ \widehat {\sigma} ^ {- 1}\tag{16}
\]

(cf. loc. cit., formula (31)).

If one identifies \(\widehat{G}\) and \(X(L^{1}(G))\) (prop. 1 of II, p.~202), the Fourier cotransformation is none other than the Gelfand transformation of the Banach algebra \(L^{1}(G)\) (I, p.~7, def. 5).

PROPOSITION 4. --- The Fourier transform and the Fourier cotransform are injective morphisms of involutive algebras from L \(^{1}\) (G) into the algebra \(\mathcal{C}_{0}(\widehat{\mathrm{G}})\) of continuous functions vanishing at infinity on \(\widehat{\mathrm{G}}\).

The Fourier cotransform is a morphism of involutive algebras from \(L^{1}(G)\) into the algebra of bounded continuous functions on \(\widehat{G}\). Since it is identified with the Gelfand transform, its image is contained in \(\mathcal{C}_{0}(\widehat{\mathrm{G}})\) (I, p.~37, prop. 5), and its kernel is the radical of \(L^{1}(G)\) (prop. 8 of I, p.~38), which is zero (cor. of prop. 22 of I, p.~126).

We shall see later (corollary of Proposition 13 of II, p.~221) that the Fourier cotransformation on \(\mathcal{M}^{1}(\mathrm{G})\) is also injective.

Remark. --- The Fourier transform on the space \(L^{1}(G)\) depends on the choice of the Haar measure dx, unlike the Fourier transform on \(\mathcal{M}^{1}(G)\). If one replaces dx by the measure \(a \cdot dx\) (with \(a > 0\)), then for any integrable function f on G, the Fourier transform of f is \(a\widehat{f}\), where \(\widehat{f}\) is the Fourier transform defined with respect to the measure dx.

Consider the stellar algebra Stell(G) of the group G (def. 9 of I, p.~125), and identify L \(^{1}\) (G) with a dense subalgebra of Stell(G) (prop. 22 of I, p.~126).

PROPOSITION 5. --- By continuity, the Fourier transform and the Fourier cotransform extend uniquely to isomorphisms of stellar algebras from Stell(G) onto \(\mathcal{C}_{0}(\widehat{\mathrm{G}})\).

The Fourier cotransform extends by continuity to a morphism of stellar algebras from Stell(G) into \(\mathcal{C}_{0}(\widehat{\mathrm{G}})\). If one identifies \(\widehat{G}\) with \(\mathrm{X}(\mathrm{Stell}(\mathrm{G}))\) (cor. 1 of II, p.~204 and prop. 1 of II, p.~202), this extension is the Gelfand transform of Stell(G). By th. 1 of I, p.~108, it is an isomorphism. The assertion concerning the Fourier transform follows from this.

We will always denote by \(\overline{\mathcal{F}}\) and \(\mathcal{F}\) the isomorphisms in Proposition 5.

COROLLARY. --- \(L\), the image of \(\mathrm{L}^{1}(\mathrm{G})\) under the Fourier transform of G, is dense in \(\mathcal{C}_{0}(\widehat{\mathrm{G}})\).

Since \(L^{1}(G)\) is dense in \(\text{Stell}(G)\), this follows from Proposition 5.

PROPOSITION 6. — Suppose that G is compact. The normalized Haar measure dx belongs to \(\mathcal{M}^{1}(\mathrm{G})\), and its Fourier transform is \(\varphi_{e}\), the characteristic function of \(\{e\}\).

Let \(\chi \in \widehat{\mathrm{G}}\). Since \(\varepsilon_y * dx = dx\) for all \(y \in \mathrm{G}\), one has

\[
\mathcal {F} (d x) (\chi) = \overline {{\langle \chi , y \rangle}} \mathcal {F} (d x) (\chi)
\]

by formula (11). If \(\chi \neq 1\) , there exists \(y \in \mathrm{G}\) such that \(\langle \chi, y \rangle \neq 1\) , hence \(\mathcal{F}(dx)(\chi) = 0\) . If \(\chi = 1\) , then \(\mathcal{F}(dx)(\chi) = \int_{\mathrm{G}} dx = 1\) since the measure \(dx\) is normalized.

